\subsection{What fresh research could afford to cost}

The bridge also gives a concrete research-budget comparison. Let $F$
be the terminal wealth multiple from fresh research before the incremental
research bill, and $U$ the multiple from reused 75/25 research, both after
the same share-trading costs. At portfolio capital $P$, the break-even
research bill is $C^*=P(1-U/F)$. For a USD3,000 bill, minimum capital is
$P^*=3000/(1-U/F)$ when $F>U$. The Top 10 results are:

\begin{minipage}{\linewidth}
\begin{center}\small\setlength{\tabcolsep}{4pt}\renewcommand{\arraystretch}{1.18}
\begin{longtable}{@{}lrrrr@{}}\toprule
Window & Raw: max bill & Raw: min capital & 75/25: max bill & 75/25: min capital\\\midrule\endhead
3$\to$4 short & 6,100 & 147,537 & 6,867 & 131,054\\
4$\to$5 short & None & None & 2,544 & 353,759\\
5$\to$6 short & None & None & None & None\\
3$\to$4 long & 37,965 & 23,706 & 94,437 & 9,530\\
\bottomrule\end{longtable}\end{center}
\end{minipage}\par

Amounts are USD; maximum bills use USD300,000 capital. ``None'' means
fresh research trails reuse even before paying for research, so no positive
bill or portfolio-size break-even exists under proportional scaling.
In the first short window, fresh raw G10 could pay up to USD6,100,
corresponding to about USD147,537 minimum capital for a USD3,000 bill.
In the longer version of that window, the budget rises to USD37,965 and
the minimum capital falls to USD23,706. The very large USD94,437 long-window
budget for fresh 75/25 inherits its Moderna concentration documented above.
These ex-post boundaries make the economic contrast tangible: a stronger
ranking policy can stretch existing research, while a deteriorating thesis
can make fresh evidence worth substantially more.
